\section{Расширение графа знаний с использованием литературных источников}

\subsection{Роль дополнительного литературного слоя}

Базовый граф знаний опирается на статистику переходов, группы культур и региональный контекст. Следующий вопрос
состоит в том, можно ли сделать пояснения более содержательными за счёт добавления внешних агрономических источников.
Для этого базовый граф знаний был расширен литературными правилами, задающими поддерживающие и предупреждающие сигналы
для отдельных переходов между культурами.

\subsection{Исходные данные для литературного расширения}

Для литературного расширения используются результаты базового графового слоя и вручную сформированная таблица
литературных правил. Первый источник задаёт исходный интерпретационный контекст, второй добавляет внешние
агрономические сигналы для отдельных переходов между культурами.

В качестве первого входа используются результаты базового графового слоя: рассмотренные примеры, локальные
пояснения по кандидатам, структура графа знаний, интерпретационные правила и статистика переходов между культурами.
Эти данные задают исходную версию графа и базовые значения графовой поддержки, с которыми затем сравнивается версия
с учётом литературных правил.

Вторым входом является компактная таблица литературных правил. Она содержит пары культур, направление сигнала, тип
интерпретации и краткое обоснование, полученное из внешних агрономических источников. В данной работе такие правила
используются как дополнительный интерпретационный слой, а не как механизм переобучения модели.

Результаты этапа представлены в виде расширенных пояснений по кандидатам, сводки по рассмотренным примерам и
сравнения базовой и расширенной версий графа. Дополнительно формируются визуализации графов с учётом литературных
правил для тех же четырёх примеров.

\subsection{Использованные внешние агрономические источники}

Основным источником для общих правил севооборота выступает стандарт USDA NRCS Conservation Crop Rotation.
Он задаёт общую логику разнообразия культур, ограничения повторов,
использования бобовых как источника азота и сокращения повторного использования пара~\cite{nrcs2014conservation_crop_rotation}.
Именно этот источник задаёт большую часть общих правил типа
предупреждения о повторе культуры, предупреждения о насыщении группы культур и поддерживающего сигнала для бобовых
культур.

Для американского контекста переходов между \texttt{corn} и \texttt{soybeans} использовано исследование
Kluger et al., в котором преимущества такой ротации оцениваются на данных по США
\cite{kluger2021rotation_corn_soy}. Эти сигналы используются как прямые поддерживающие правила
\texttt{corn} $\rightarrow$ \texttt{soybeans} и \texttt{soybeans} $\rightarrow$ \texttt{corn}. Для перехода
\texttt{corn} $\rightarrow$ \texttt{wheat} используется осторожный предупреждающий сигнал из материалов Oklahoma
State University Extension: остатки кукурузы могут повышать риск фузариоза колоса (Fusarium Head Blight) для
последующей пшеницы, но такой сигнал трактуется как предупреждающий, а не как запрет или универсальное
агрономическое предписание \cite{osu2017corn_residue_wheat_fhb}.

Дополнительно используется методическая оговорка из современной работы Kluger et al. о том, что эффекты
предшественников зависят от последовательности культур, погодных условий и регионального контекста
\cite{kluger2025precrop_effects}. В данном этапе эта публикация служит источником контекстной интерпретации, а не
прямым источником переходных правил.

\subsection{Способ интеграции литературных правил в граф}

Литературные сигналы не подмешиваются к переходам из данных. Для каждого
сработавшего литературного правила в графе добавляется отдельный узел интерпретации типа \texttt{support} или
\texttt{warning}, связанный с кандидатом из ранжированного списка. Тем самым в расширенном графе сохраняется
различие между двумя типами сигналов:

\begin{itemize}
\tightlist
\item
  сигналами, наблюдаемыми в данных проекта;
\item
  сигналами, пришедшими из внешних агрономических источников.
\end{itemize}

Часть правил задаётся на прямом уровне для конкретной пары культур, а часть --- на уровне
\texttt{group\_approx}, когда источник говорит о семействе или агрономической группе, а не об отдельном агрегированном
классе проекта. Для точных повторов культуры применяется более специфичное предупреждение о повторе. В этом случае
предупреждение о насыщении той же группы культур не добавляется повторно, чтобы не учитывать один и тот же фактор
дважды. Для неполных повторов внутри одной группы такое предупреждение сохраняется.

Для класса \texttt{forage\_hay} предупреждение о повторе задано как слабое, поскольку агрегированный класс может
включать многолетние или кормовые звенья, роль которых отличается от повторного выращивания однолетних пропашных
культур.

На уровне отдельных кандидатов это приводит к пересчёту \texttt{graph\_support\_score}. Новый показатель суммирует
более широкий набор сигналов поддержки и предупреждения. Тем самым литературные правила изменяют интерпретационную
поддержку отдельных кандидатов.

\subsection{Результаты на тех же примерах}

Литературное расширение было применено к тем же четырём примерам, что и в базовом интерпретационном анализе:
\texttt{case\_high\_confidence}, \texttt{case\_medium\_confidence}, \texttt{case\_low\_confidence} и
\texttt{case\_difficult\_class}. Это позволяет сравнивать один и тот же набор ситуаций до и после добавления
литературного слоя.

На выбранных четырёх примерах литературные правила не изменили итоговую категорию согласованности, но изменили
интерпретационную поддержку отдельных кандидатов. В \texttt{case\_high\_confidence} сохранился статус
\texttt{strong}, а для трёх остальных примеров сохранился статус \texttt{weak}. При этом внутри отдельных
ранжированных списков изменились значения графовой поддержки и набор сработавших интерпретационных правил.

Сводное сравнение по тем же примерам приведено в таблице~\ref{tab:literature-effect}. Она показывает, что
литературные правила влияют прежде всего на интерпретационную поддержку отдельных кандидатов, а не на итоговую
категорию согласованности.

\begin{table}[H]
\centering
\scriptsize
\caption{Изменение графовой поддержки после добавления литературных правил}
\label{tab:literature-effect}
\setlength{\tabcolsep}{2pt}
\begin{tabular}{@{}>{\raggedright\arraybackslash}p{0.13\textwidth}
                >{\raggedright\arraybackslash}p{0.12\textwidth}
                r
                r
                r
                >{\raggedright\arraybackslash}p{0.18\textwidth}
                >{\raggedright\arraybackslash}p{0.14\textwidth}@{}}
\toprule
Пример & \makecell[l]{Первая рекомендация\\модели} & \makecell[r]{Базовый\\$gs$} & \makecell[r]{$gs$ с\\литературными\\правилами} & \makecell[r]{$\Delta gs$\\первой\\рекомендации}
    & \makecell[l]{Кандидат с\\наибольшим $gs$} & Согласованность \\
\midrule
Высокая уверенность
  & \texttt{forage\_hay}
  & 0.661
  & 0.591
  & $-0.070$
  & \texttt{forage\_hay} (0.591)
  & сильная $\rightarrow$ сильная \\
Средняя уверенность
  & \texttt{corn}
  & 0.376
  & 0.236
  & $-0.140$
  & \texttt{soybeans} (0.598)
  & слабая $\rightarrow$ слабая \\
Низкая уверенность
  & \texttt{legumes}
  & 0.011
  & 0.071
  & $+0.060$
  & \texttt{soybeans} (0.698)
  & слабая $\rightarrow$ слабая \\
Сложный класс
  & \texttt{corn}
  & 0.220
  & 0.220
  & $\phantom{+}0.000$
  & \texttt{soybeans} (0.386)
  & слабая $\rightarrow$ слабая \\
\bottomrule
\end{tabular}

\vspace{0.4em}
\footnotesize $gs$ --- \texttt{graph\_support\_score}; $\Delta gs$ --- разность между значением с учетом литературных правил
и базовым значением для первой рекомендации модели.
\normalsize
\end{table}


В данном разделе рассматриваются не метрики предиктивного качества модели, а изменения интерпретационных показателей
графа для уже выбранных примеров. После фильтрации было зафиксировано 9 срабатываний литературных правил; при этом для 3 строк с
точным повтором культуры предупреждение о насыщении той же группы было пропущено, поскольку более специфичный сигнал
повтора уже был учтён. На уровне рассмотренных примеров это особенно заметно для \texttt{forage\_hay}, где повтор
интерпретируется как слабое предупреждение, и для \texttt{corn}, где повтор культуры ослабляет поддержку первой
рекомендации, тогда как переход к \texttt{soybeans} получает дополнительную поддержку.

Если рассмотреть усреднённую динамику по всем кандидатам, средняя разница между расширенным и базовым значениями
графовой поддержки оказывается небольшой, но положительной и составляет примерно 0.008. Для первых рекомендаций
модели средняя дельта отрицательна и составляет около $-0.020$, тогда как для остальных кандидатов она положительна
и составляет около 0.023. На уровне отдельных кандидатов изменения оказались разнонаправленными: часть кандидатов
получила дополнительную поддержку, часть --- предупреждающие сигналы, а часть практически не изменилась.

\subsection{Интерпретация результатов}

Основной эффект литературного слоя проявляется не в изменении ранжирования, а в уточнении пояснений к отдельным кандидатам.
Литературный слой делает объяснения на уровне
отдельных кандидатов более содержательными, позволяет явно фиксировать предупреждающие сигналы для повторных или
нежелательных переходов и усиливает некоторые агрономически правдоподобные альтернативы в ранжированном списке.
Это особенно полезно в спорных случаях, где пользователю важно видеть не один ответ, а набор конкурирующих
интерпретируемых альтернатив.

При интерпретации этих результатов нужно учитывать несколько особенностей. Во-первых, сам набор литературных правил формируется вручную и
остаётся относительно компактным. Во-вторых, часть правил построена на уровне \texttt{group\_approx}, поскольку
агрегированные классы проекта не всегда имеют прямое соответствие одной конкретной культуре из источника.
В-третьих, для класса \texttt{cotton} в литературном слое используются только общие предупреждения, связанные с
повтором культуры и недостаточным разнообразием. Наконец, анализ выполнялся на тех же четырёх примерах, а не на
полном тестовом наборе, поэтому он характеризует прежде всего качество интерпретации.

Таким образом, литературное расширение повышает содержательность пояснений в рамках дополнительного
интерпретационного слоя. CatBoost формирует ранжированный список кандидатов, слой оценки уверенности характеризует
надёжность результата, а граф знаний добавляет предметный контекст для интерпретации отдельных вариантов. Внешние
агрономические источники в этом контуре используются как дополнительная аргументация для уже построенных разборов
примеров.
